\chapter[Introdução]{Introdução}
\label{cap:introducao}

Parentesco é a relação que une duas ou mais pessoas por vínculos genéticos de descendência ou ascendência, e também por vínculos sociais como adoção e convivência. Este trabalho trata especificamente da verificação visual de parentesco consanguíneo, ou seja, da decisão sobre se duas pessoas têm vínculo biológico a partir apenas de imagens faciais. Vínculos sociais sem base genética ficam fora do escopo, porque o sinal explorado pela visão computacional é a semelhança de traços herdados.

Verificação de parentesco tem sido objeto de estudo constante no campo da psicologia. Pessoas com vínculos genéticos podem possuir certo nível de semelhanças físicas entre si, e \citeonline{hogben1932genetic} denominou essas similaridades na estrutura facial como traços familiares. A semelhança facial é considerada uma das mais comuns referências para indicar relações de parentesco \cite{debruine2008social}. A hipótese de que a similaridade entre faces poderia ser uma pista para sugerir parentesco foi primeiramente formulada por \citeonline{daly1982whom}. Maloney e Martello examinaram a relação entre semelhança das estruturas faciais e a existência ou não de relações de consanguinidade entre irmãos e concluíram que observadores tendem a procurar por similaridades na verificação de parentesco entre crianças \cite{maloney2006kin}.

Em 1984, \citeonline{porter1984recognition} registrou que pessoas conseguiam parear fotos de mães com seus respectivos filhos sem conhecê-los, enquanto as próprias mães reconheciam seus bebês poucas horas após o nascimento. Isso reforça a tese de que há de fato um sinal facial explorável entre parentes próximos.

A verificação automática desse vínculo a partir de fotos de rosto é um campo relativamente novo e difícil em visão computacional. O sistema decide, diante de um par de imagens, se as duas pessoas têm ou não relação de parentesco, examinando apenas estruturas faciais. O problema é particularmente difícil por duas razões opostas, porque parentes próximos podem apresentar traços bastante distintos, como em grupos de irmãos nos quais um herda mais da mãe e outro mais do pai, e, ao mesmo tempo, pessoas sem vínculo biológico podem ter rostos muito parecidos. Em relação ao reconhecimento facial comum a tarefa é mais fina, porque não basta separar identidades, é preciso capturar padrões de herança em imagens que variam em idade, expressão, iluminação e pose.

O reconhecimento de relações consanguíneas por imagem tem usos práticos concretos como auxiliar investigações de tráfico humano, identificação de pessoas com dificuldade de comunicação (crianças e idosos desaparecidos, por exemplo), montagem de árvores genealógicas e redução do espaço de busca em bancos de faces em larga escala. A verificação de parentesco por face também pode servir como apoio forense à identificação de crianças quando a documentação está ausente ou foi adulterada.

Os conflitos armados em curso (Ucrânia, Gaza, Sudão, Mianmar, Sahel, República Democrática do Congo e Iêmen) colocam o mesmo problema em outra escala. Só na Ucrânia, investigações conduzidas pela Europol apontam cerca de 19.500 crianças transferidas à força para territórios controlados pela Rússia e pela Belarus, e o rastreamento depende hoje principalmente de dados genéticos \cite{europol_ucrania_2026, rtp_dna_ucrania_2026}. Análise facial automática não substitui DNA, mas pode servir como triagem em postos de identificação e em programas de family tracing conduzidos por organizações como CICV, ACNUR e UNICEF.

As tentativas de construir um modelo computacional de verificação visual de parentesco começaram em 2010 e estão descritas em \citeonline{fang2010towards}, que combina atributos faciais extraídos a partir de estrutura pictórica \cite{felzenszwalb2005pictorial} com classificação por Support Vector Machine, alcançando cerca de 68{,}6\% de acurácia em um conjunto próprio de 150 pares. Em 2018, \citeonline{robinson2018visual} disponibilizou Families in the Wild, hoje o maior banco público de parentesco facial, e com ele se difundiu o uso de redes convolucionais pré-treinadas em reconhecimento de face, como VGG-Face \cite{parkhi2015deep} e ResNet \cite{wen2016discriminative}. Desde então a literatura migrou para metric learning com perdas contrastivas, mecanismos de atenção e Vision Transformers, e os melhores resultados recentes no FIW ficam próximos de 92\% de acurácia.

Apesar desse avanço, comparar os métodos da área com rigor continua trabalhoso. A literatura usa protocolos heterogêneos, com thresholds frequentemente ajustados sobre o próprio conjunto de teste, métricas diferentes entre artigos e poucos estudos cruzando mais de um banco sob a mesma receita. Ao mesmo tempo, modelos multimodais generalistas (VLMs) ganharam popularidade e começam a sugerir que talvez nem seja preciso treinar um classificador específico para parentesco, podendo a tarefa ser resolvida por inferência direta sobre um modelo multimodal sem treinamento de domínio.

\section{Problema de pesquisa}

A literatura recente em verificação visual de parentesco facial reúne arquiteturas heterogêneas avaliadas sob protocolos não comparáveis, e os ganhos relatados frequentemente dependem de escolhas de threshold, amostragem e seleção de checkpoint que variam entre trabalhos. Em paralelo, modelos multimodais generalistas em regime zero-shot têm sido apresentados como alternativa de baixo custo de implantação, mas faltam medições rigorosas que situem o desempenho desses modelos contra modelos supervisionados especializados, considerando que as duas classes atacam tarefas distintas (verificação binária para os supervisionados e classificação multiclasse das relações para os VLMs).

Diante disso, o problema de pesquisa deste trabalho se desdobra em duas perguntas complementares. A primeira é a comparação principal:

\begin{quote}
\noindent Em que medida uma arquitetura supervisionada especializada para parentesco facial, combinando backbone de reconhecimento facial, descongelamento parcial e tokens regionais anatômicos, supera um baseline sem treinamento e arquiteturas supervisionadas alternativas no benchmark Families in the Wild, sob protocolo único de verificação binária?
\end{quote}

A segunda é uma pergunta secundária, de caráter exploratório:

\begin{quote}
\noindent Qual é o desempenho atual de modelos multimodais generalistas (VLMs) em classificação fina de relações de parentesco no FIW, em regime zero-shot e em adaptação por prompt, quando comparados qualitativamente aos modelos supervisionados em uma tarefa relacionada porém distinta?
\end{quote}

\section{Objetivo geral}

Comparar de forma controlada cinco modelos representativos para verificação binária visual de parentesco facial, sob um único protocolo experimental no benchmark Families in the Wild, e medir o ganho efetivo de uma proposta arquitetônica baseada em backbone facial especializado, descongelamento parcial e tokens regionais anatômicos. Em paralelo e como estudo exploratório complementar, mapear o desempenho de dois modelos multimodais generalistas em classificação multiclasse de relações de parentesco no mesmo banco.

\section{Objetivos específicos}

\begin{enumerate}
  \item Implementar um baseline mínimo de verificação de parentesco com extrator facial AdaFace pré-treinado e similaridade cosseno, sem qualquer treinamento específico de domínio.
  \item Reproduzir e avaliar três famílias de modelo da literatura recente em regime de verificação binária, contemplando ViT com cross-attention bidirecional, DINOv2 adaptado por LoRA com cross-attention diferencial e modelo retrieval-augmented sobre galeria de pares positivos.
  \item Implementar e avaliar uma proposta autoral para o Modelo 12, que combina AdaFace IR-101 com descongelamento parcial do último estágio, cinco tokens regionais anatômicos, cross-region attention bidirecional, gate sigmoide, cabeça auxiliar de relação e inferência simétrica entre as duas ordens do par.
  \item Conduzir uma avaliação exploratória de dois modelos multimodais generalistas, Codex e Claude Sonnet, em classificação multiclasse das onze relações do FIW, em regime zero-shot e em adaptação ao domínio por prompt, sem pretensão de comparação direta às métricas binárias dos modelos supervisionados.
  \item Reportar e discutir métricas robustas a operating point específico, em particular ROC AUC e Average Precision, e métricas calculadas em pontos operacionais fixados por taxa de falso positivo (TAR@FAR=0,1, TAR@FAR=0,01 e TAR@FAR=0,001), além de análise por tipo de relação, identificando limitações metodológicas e éticas do uso pretendido.
\end{enumerate}

\section{Justificativa}

A relevância do trabalho tem duas frentes. Do ponto de vista científico, a área de verificação visual de parentesco facial carece de comparações controladas entre famílias de método sob o mesmo protocolo, e a entrada recente de modelos multimodais generalistas adiciona uma dimensão pouco mapeada na literatura. Do ponto de vista aplicado, a tarefa tem usos potenciais em apoio forense, reunificação familiar em contextos humanitários e identificação assistida em programas de family tracing. A delimitação ética desse uso, no entanto, é estrita. O sistema descrito aqui é uma ferramenta de triagem ou apoio à decisão humana qualificada e nunca substitui prova genética ou avaliação especializada. Em particular, falsos positivos têm custo alto em contextos forenses e em populações vulneráveis (crianças desaparecidas, refugiados, civis em zonas de conflito), e a tecnologia deve operar sob controles institucionais, consentimento adequado, rastreabilidade da decisão e direito à contestação por parte das pessoas afetadas. Diferenças demográficas conhecidas em modelos de reconhecimento facial \cite{buolamwini2018gender} sugerem que viés por grupo é um risco real também em verificação de parentesco, e a ausência de medições específicas de fairness nesta versão do trabalho é declarada como limitação.

\section{Contribuições}

As contribuições principais deste trabalho são quatro. A primeira é um protocolo comparativo aplicado de forma uniforme aos cinco modelos supervisionados de verificação binária, com threshold escolhido em validação, métricas robustas a operating point e análise por tipo de relação. A segunda é a proposta arquitetônica autoral do Modelo 12, descrita no Capítulo~\ref{cap:Metodologia}, que combina backbone facial especializado, descongelamento parcial do último estágio, estrutura regional anatômica com cross-region attention, gate sigmoide, cabeça auxiliar de relação e forward simétrico para reduzir atalhos dependentes da ordem do par. A terceira é uma avaliação exploratória de dois VLMs em parentesco facial, em classificação multiclasse das onze relações do FIW, com a ressalva explícita de que se trata de tarefa diferente da verificação binária dos supervisionados. A quarta é uma análise consolidada de erros por tipo de relação, com atenção particular às relações de avós e netos, classes historicamente mais difíceis na literatura do FIW.

\section{Delimitações}

Este trabalho restringe-se à verificação visual de parentesco consanguíneo a partir de imagens faciais. Vínculos sociais sem base genética e parentesco identificado a partir de DNA, documentos ou registros oficiais ficam fora do escopo. A avaliação principal usa o benchmark Families in the Wild com o split Track-I oficial do RFIW, com validação cruzada de cinco dobras restrita ao KinFaceW-I como leitura complementar para o Modelo 02. O trabalho não pretende propor um sistema operacional de identificação, e não foram conduzidas medições controladas de viés demográfico ou avaliação cross-dataset entre KinFaceW-I e FIW.

\section{Organização do trabalho}

O restante deste trabalho está organizado da seguinte maneira. O Capítulo~\ref{cap:trabalhos-relacionados} situa o trabalho na literatura recente, do primeiro modelo computacional para a tarefa aos modelos com Vision Transformer e atenção cruzada. O Capítulo~\ref{cap:fundamentacao-teorica} apresenta os conceitos necessários para compreender as escolhas metodológicas. O Capítulo~\ref{cap:Metodologia} descreve datasets, modelos, treinamento, protocolo de avaliação e baselines com VLMs. O Capítulo~\ref{cap:resultados} reporta os resultados experimentais, e o Capítulo~\ref{cap:conclusao} discute os achados, limitações e trabalhos futuros.
